\chapter{Formal logic as a working tool}

\begin{orient}
\textbf{What this chapter is for.} Almost every logic puzzle arrives written in English, and English is ambiguous in ways that matter: ``if'' sometimes means ``if and only if'', ``or'' sometimes excludes, ``not both'' is routinely read as ``neither''. The formal apparatus of propositional and predicate logic earns its place here for one reason only, which is that it converts those sentences into objects you can compute with. A truth table is a finite computation. A contrapositive is a rewrite that preserves meaning exactly. De Morgan's laws and the quantifier negation rules are a mechanical procedure for finding out what it would take to refute a claim, which is very often the only usable handle a problem offers. This chapter covers exhaustion over assignments and the point at which it stops being practical, the four forms a conditional can take and which of them you actually own, negation normal form for compound and quantified statements, the material conditional together with the two traps it sets, and the discipline of separating necessary from sufficient when reading a problem statement. By the end you should be able to take a paragraph of puzzle English, write down the exact propositional content, and know before starting whether the question is one of satisfiability, of entailment, or of optimality.
\end{orient}

\section{Exhaustion, and where it stops}

The cheapest thing you can do with a set of propositional constraints is to try every assignment. It is worth being precise about what that computation returns, because three different questions are routinely asked of the same table and they have three different answers.

Fix atomic propositions $p_1,\ldots,p_n$ and constraints $C_1,\ldots,C_m$, and write $\varphi=C_1\wedge\cdots\wedge C_m$. The table of $2^n$ rows partitions the assignments into those satisfying $\varphi$ (call them the \emph{survivors}) and those falsifying it. \emph{Satisfiability} asks whether at least one survivor exists. \emph{Entailment}, written $\varphi\models\psi$, asks whether every survivor also satisfies $\psi$; this is not a question about any single row. \emph{Determination} asks whether all survivors agree on the value of some particular atom, which is the question a puzzle means when it asks ``who is the knight?''. A puzzle can be satisfiable and yet determine nothing, and the difference between those two outcomes is the difference between a well-posed puzzle and a badly posed one.

\tech{Truth-table exhaustion}{easy}
\cue A finite number of atomic propositions, in practice $n\leq4$ by hand and $n\leq20$ by machine, tied together by compound constraints, and a question of the form ``is this consistent'', ``does this follow'', or ``which assignments survive''. Also any puzzle whose entire content is the joint satisfiability of a list of English sentences.
\method Translate each sentence into a formula over the atoms, enumerate all $2^n$ assignments, and keep the rows on which every constraint evaluates to true. Read the answer off the surviving set: nonempty means consistent, a single row means fully determined, and $\varphi\models\psi$ holds exactly when no survivor falsifies $\psi$. The one operation that is never legitimate is inspecting a favourable row and stopping.

\begin{worked}
Three relays $p,q,r$ are wired so that $p\to q$, $q\to r$, and $\lnot(p\wedge r)$. Which configurations are possible, and does it follow that $p$ is off?

Take the cases on $p$. If $p$ is true then $q$ is true by the first constraint and $r$ is true by the second, so $p\wedge r$ holds and the third constraint fails. Hence $p$ is false in every survivor, and with $p$ false the first constraint is automatically satisfied, leaving $q\to r$ as the only live requirement. The survivors are therefore
\[(p,q,r)=(F,F,F),\quad (F,F,T),\quad (F,T,T),\]
three rows out of eight, confirmed by enumerating all eight assignments in \texttt{python3}. So $\varphi\models\lnot p$: the relay $p$ is off in every consistent configuration. Note what does \emph{not} follow. Two survivors have $q$ false and one has $q$ true, so $\varphi\not\models\lnot q$ and $\varphi\not\models q$: the state of $q$ is undetermined, and any answer that names it is wrong.
\end{worked}

\begin{trap}
Two failures, and they are opposite. The first is exhibiting one consistent row and reporting it as the answer, which establishes satisfiability and nothing more. The second is finding one row on which the conclusion happens to hold and calling that entailment; entailment is a statement about every survivor, so it is refuted by a single survivor and confirmed by none.
\end{trap}

\begin{seealsob}De Morgan normalisation; the Boolean encoding of a truthful or lying speaker; exhaustive casework with a discipline.\end{seealsob}

The table is a complete decision procedure, which is exactly why it is the wrong tool as soon as the problem has any structure. Doubling $n$ squares the work, so $n=30$ is already out of reach for a table you write out, and even at $n=6$ the sixty-four rows are a transcription exercise in which arithmetic slips are more likely than insight. Every remaining technique in this chapter is a way of not building the table: a rewrite that halves the case count, a normalisation that makes the branching visible, or a reading discipline that stops you from formalising the wrong sentence in the first place. The single most valuable of those rewrites is the one that turns a conditional around.

\section{Rewriting a conditional}

A conditional $A\to B$ has three relatives, and puzzles are built on the fact that only one of them is equivalent to it. The equivalence is not a curiosity. It is the standard move for turning an unusable hypothesis into a usable one, because a problem is almost always engineered so that exactly one of $A$ and $B$ is the concrete side.

\begin{center}
\footnotesize
\setlength{\tabcolsep}{7pt}
\renewcommand{\arraystretch}{1.25}
\begin{tabular}{@{}cc|cccc@{}}
\toprule
\mh{$A$} & \mh{$B$} & \mh{$A\to B$} & \mh{$B\to A$} & \mh{$\lnot A\to\lnot B$} & \mh{$\lnot B\to\lnot A$} \\
 & & \mh{original} & \mh{converse} & \mh{inverse} & \mh{contrapositive} \\
\midrule
T & T & T & T & T & T \\
T & F & F & T & T & F \\
F & T & T & F & F & T \\
F & F & T & T & T & T \\
\bottomrule
\end{tabular}
\end{center}

Read the table by columns, not by rows. The first and fourth columns are identical, which is the whole content of the contrapositive rule. The second and third are identical to each other and differ from the first in the two middle rows, which is why a puzzle that hands you the converse has handed you nothing about the original. And the first column is false in exactly one row, a fact so consequential that it gets its own technique later in this chapter.

\tech{Contrapositive, converse, inverse}{easy}
\cue A hypothesis of the form ``if $A$ then $B$'', and a problem that gives you information about $B$ or about $\lnot B$ rather than about $A$. The moment a puzzle offers you the consequent, stop and ask which direction you actually own.
\method Only $\lnot B\to\lnot A$ is equivalent to $A\to B$; the converse $B\to A$ and the inverse $\lnot A\to\lnot B$ are equivalent to each other and to neither. So a rule ``$A$ implies $B$'' is tested by looking at things that are $A$ and at things that are \emph{not} $B$, and by nothing else. In symbols,
\[A\to B\ \equiv\ \lnot B\to\lnot A\ \equiv\ \lnot A\vee B,\]
and the last form is the one to use when you want to case-split rather than chain.

\begin{worked}
A warehouse rule states: every crate bearing a red seal contains glass. Four crates sit on the dock. You can see that the first has a red seal, that the second has a blue seal, that the third contains glass, and that the fourth contains steel; in each case the other fact is hidden. Which crates must be opened or unsealed to test the rule?

Write $A$ for ``red seal'' and $B$ for ``contains glass''; the rule is $A\to B$. Crate one shows $A$, so it tests the rule directly: if it holds steel the rule is dead. Crate four shows $\lnot B$, so it tests the contrapositive $\lnot B\to\lnot A$: if its seal is red the rule is dead. Crate two shows $\lnot A$ and crate three shows $B$, and by the table above neither row can falsify $A\to B$; whatever they hide, the rule survives. Answer: crates one and four, and only those.

The instructive part is why crate three is so tempting. Opening it can only confirm the rule, and confirmation is not a test. A rule of the form $A\to B$ is refutable in exactly one way, by an $A$ that is not a $B$, so the crates worth touching are precisely those that could be that counterexample.
\end{worked}

\begin{trap}
Inspecting the case that could confirm the rule instead of the case that could refute it. This is the most reliable single failure in conditional reasoning, and it survives knowing the truth table, because the pull is toward the row that makes the rule look good. The mechanical guard: before choosing what to examine, write the unique falsifying row $A$ true and $B$ false, and examine only what could occupy it.
\end{trap}

\begin{seealsob}Necessary versus sufficient; proof by contrapositive; the material conditional and vacuous truth.\end{seealsob}

\begin{keynote}[The one row that matters]
Everything about the conditional follows from the observation that $A\to B$ is false in exactly one of its four rows. To \emph{refute} a conditional, produce that row. To \emph{use} a conditional, note that it forbids that row and nothing else. To \emph{prove} a conditional, rule that row out, which you may do by assuming $A$ and deriving $B$ or by assuming $\lnot B$ and deriving $\lnot A$; those are the two ways of walking into the same forbidden cell. A great many arguments that feel like different methods are this one observation approached from different sides.
\end{keynote}

\section{Pushing a negation inward}

The rewrites of the previous section handle a negation sitting outside a conditional. Negations also sit outside conjunctions, disjunctions and quantifiers, and there the rule is uniform: push the negation in one connective at a time, flipping each connective as it passes, until every negation sits on an atom. The resulting form is called negation normal form, and its virtue is that the case structure of the problem becomes visible on the page instead of remaining hidden inside a $\lnot$.

\tech{De Morgan normalisation}{easy}
\cue A negated compound, in any of its English disguises: ``it is not the case that both'', ``none of them'', ``at most one of these holds'', ``neither''. Also any moment when you want to case-split but the statement offers you no visible disjunction to split on.
\method Apply
\[\lnot(A\wedge B)\equiv\lnot A\vee\lnot B,\qquad \lnot(A\vee B)\equiv\lnot A\wedge\lnot B,\]
repeatedly, together with $\lnot(A\to B)\equiv A\wedge\lnot B$ and $\lnot\lnot A\equiv A$, until no negation has a compound in its scope. Then read the branching directly: a top-level $\vee$ is a case split, a top-level $\wedge$ is a list of facts you may all use at once.

\begin{worked}
A label on a box reads: ``it is not true both that the prize is in this box and that this label is honest.'' Normalise.

Let $P$ be ``the prize is in this box'' and $H$ ``this label is honest''. The label asserts $\lnot(P\wedge H)$, which normalises to $\lnot P\vee\lnot H$: either the prize is elsewhere or the label lies, possibly both. That is two branches, and it is worth seeing that the sentence has not been weakened by the rewrite, only made usable. Now add the self-referential constraint from the next chapter's encoding, $H\leftrightarrow\lnot(P\wedge H)$. If $H$ is true then $\lnot(P\wedge H)$ is true, so $P\wedge H$ is false, so $P$ is false. If $H$ is false then $\lnot(P\wedge H)$ is false, so $P\wedge H$ is true, so $H$ is true, a contradiction. Answer: the label is honest and the prize is elsewhere, the unique surviving row of four.
\end{worked}

\begin{worked}[The counting version]
``At most one of $A_1,A_2,A_3$ holds'' is a negated existential over pairs: $\lnot\exists i<j\,(A_i\wedge A_j)$, which normalises to $\bigwedge_{i<j}(\lnot A_i\vee\lnot A_j)$, three clauses. ``At least two hold'' is its negation, $\bigvee_{i<j}(A_i\wedge A_j)$, three cases. ``Exactly one holds'' is the conjunction of ``at most one'' with $A_1\vee A_2\vee A_3$. Written this way the three phrases become three different pieces of syntax and stop being confusable, which is the entire practical point of normalising.
\end{worked}

\begin{trap}
Reading ``not both'' as ``neither''. The De Morgan output is an inclusive disjunction, and that inclusive $\vee$ is exactly what makes the puzzle branch instead of collapse. The mirror error is reading ``neither'' as ``not both'', which throws away a constraint: $\lnot A\wedge\lnot B$ is strictly stronger than $\lnot A\vee\lnot B$.
\end{trap}

\begin{seealsob}Quantifier negation; truth-table exhaustion; existential claims about the speakers themselves.\end{seealsob}

The same procedure runs over quantifiers, and there it does more work, because a quantified statement usually cannot be attacked head-on while its negation often has a one-object proof. Knowing exactly what the negation says is therefore not bookkeeping; it is the step that decides whether you are looking for a construction or for an argument over an infinite domain.

\tech{Quantifier negation}{medium}
\cue Any statement built from ``all'', ``some'', ``no'', ``at least one'', ``exactly one'', which you need to refute, or which appears inside a negation, or whose failure is the constructive lever the problem wants.
\method Push the negation inward one quantifier at a time, flipping each as it passes and leaving the order untouched:
\[\lnot\forall x\,P(x)\equiv\exists x\,\lnot P(x),\qquad \lnot\exists x\,P(x)\equiv\forall x\,\lnot P(x).\]
For nested quantifiers apply this outside-in, so that $\lnot\forall x\,\exists y\,R(x,y)\equiv\exists x\,\forall y\,\lnot R(x,y)$. The asymmetry is the payoff: refuting a universal costs exactly one witness, whereas refuting an existential costs a proof over the whole domain. Read the negated form first and let it tell you which of those two jobs you have.

\begin{worked}
Claim: in every red-blue colouring of the edges of the complete graph on five vertices there is a triangle with all three edges the same colour. Negate it. The negation is: there exists a colouring in which every triangle has two colours among its edges. That is an existential over a finite set, so one explicit colouring settles the matter.

Label the vertices $0,1,2,3,4$ and colour the edge $\{i,j\}$ red when $j-i\equiv\pm1\pmod5$ and blue otherwise. The red edges form a pentagon and the blue edges form the pentagram; neither is a graph containing a triangle, since both are five-cycles. So no triangle is monochromatic and the claim is false. Exhaustive search over all $2^{10}=1024$ colourings in \texttt{python3} confirms this and returns $12$ such colourings in all.

Now run the same negation on six vertices. There the negation asserts a colouring of $\binom{6}{2}=15$ edges with no monochromatic triangle, and exhaustive search over all $2^{15}=32768$ colourings returns none. So on six vertices the universal claim is true, and the negation has told us that no witness exists rather than exhibiting one. Answer: false for five vertices, true for six.
\end{worked}

\begin{trap}
Two errors, both silent. The first is negating $\forall x\,P(x)$ to $\forall x\,\lnot P(x)$, the strong denial: ``not everything is prime'' is not ``everything is composite''. The second is reordering quantifiers while negating. The statements $\forall x\,\exists y\,R(x,y)$ and $\exists y\,\forall x\,R(x,y)$ are genuinely different, the second being strictly stronger, and a negation that shuffles them is answering a different question.
\end{trap}

\begin{seealsob}De Morgan normalisation; the extremal principle; proof by contradiction.\end{seealsob}

\begin{keynote}[Order of quantifiers, in one example]
Over the positive integers, $\forall m\,\exists n\,(n>m)$ is true, since $n=m+1$ works for each $m$ separately. Swapping the quantifiers gives $\exists n\,\forall m\,(n>m)$, which asserts a single integer exceeding every integer and is false. Nothing changed but the order, and the truth value flipped. In puzzle English the distinction hides inside the word ``a'': ``for every lock there is a key'' and ``there is a key for every lock'' are the two readings, and the second is the one a careless problem statement means by accident. When a problem says ``there is'', ask immediately whether the object is allowed to depend on the earlier variables.
\end{keynote}

\section{What ``if'' actually asserts}

The connective $\to$ was defined above by a column of a truth table, and that definition has two consequences that regularly surprise readers. The first is that a conditional with a false antecedent is true, no matter what the consequent says. The second is that English ``if'' frequently means something stronger than $\to$, so that formalising a sentence as $A\to B$ can lose information the problem intended to give you. Both consequences are exploitable, and both are traps.

\tech{The material conditional and vacuous truth}{medium}
\cue A conditional whose antecedent might be false: a speaker's conditional claim, a rule quantified over a set that may turn out to be empty, or a hypothesis of the form ``if the first box is empty then \ldots''.
\method Take $A\to B$ to mean $\lnot A\vee B$, true in every row but $A$ true with $B$ false. Two working consequences. A conditional is established for free once you know its antecedent fails, which is \emph{vacuous truth}: a rule such as ``every element of $S$ is red'' is true when $S$ is empty. And an assertion of $A\to B$ by someone whose truthfulness is itself in question can force $A$ to be true, because the assertion is cheap when $A$ fails and expensive when it holds.

\begin{worked}
On an island where each inhabitant is either a knight, who always speaks truly, or a knave, who always speaks falsely, one inhabitant says: ``if I am a knight, then this box holds gold.'' What follows?

Write $k$ for ``the speaker is a knight'' and $G$ for ``the box holds gold''. The speaker's assertion is $S=(k\to G)$, and the island rule is that the speaker is a knight exactly when the assertion is true, so $k\leftrightarrow S$. Suppose $k$ is false. Then $S=(F\to G)=T$, so $k\leftrightarrow S$ reads $F\leftrightarrow T$, which fails. So $k$ is true, hence $S$ is true, hence $G=(k\to G)$ with $k$ true gives $G$ true. Answer: the speaker is a knight and the box holds gold, the unique consistent row among the four, verified by enumerating all four assignments in \texttt{python3}.

Change one word and everything changes. If the speaker instead says ``if I am a knave, then this box holds gold'', the assertion is $\lnot k\to G$, and enumeration returns three consistent rows: knight with gold, knight without gold, and knave without gold. Nothing at all is determined. The asymmetry is pure vacuous truth: ``if I am a knight'' has an antecedent that a knight cannot escape, whereas ``if I am a knave'' has an antecedent that a knight discharges for free.
\end{worked}

\begin{trap}
Importing relevance into $\to$. The sentence ``if the moon is cheese then $2+2=5$'' is true, and no amount of discomfort changes that; there is no causal content in the connective at all. The mirror trap is the second one this section promised: English ``if'' is very often biconditional. ``You may enter if you have a ticket'' is normally meant as ``if and only if''. When a puzzle says ``if'', decide from the context whether the converse was also intended, and if the problem turns on the difference, say which reading you took.
\end{trap}

\begin{seealsob}Truth-table exhaustion; the Boolean encoding of a truthful or lying speaker; necessary versus sufficient.\end{seealsob}

\begin{keynote}[Vacuous truth is not a loophole]
``Every purple crate on this dock contains gold'' is true when no purple crate is present, and this is a consequence of the definition rather than a lawyer's reading of it. The reason to define $\to$ this way is that the alternative destroys induction and quantification: we want $\forall n\,(n>100\to P(n))$ to be checkable by looking only at $n>100$, and we want a statement about the empty set to be harmless rather than meaningless. In puzzles the practical consequences are two. A universal statement over a possibly empty range gives you nothing until you have shown the range is nonempty. And a claim you cannot refute may be unrefutable because it is vacuous, in which case the useful move is to attack the emptiness rather than the claim.
\end{keynote}

Vacuous truth is what happens when a conditional's antecedent is empty. The remaining reading discipline concerns which side of a conditional a piece of English has landed on, and it is the one place in this chapter where the cost of a mistake is an entire wrong answer rather than a lost step.

\tech{Necessary versus sufficient}{medium}
\cue The words ``only if'', ``provided that'', ``it suffices that'', ``one must'', ``no more than''. Also, and this is the important case, any problem asking for a \emph{minimum} number of tests, questions, weighings or moves, where a bound and a construction are two separate obligations.
\method Fix the dictionary and use it mechanically. ``$A$ only if $B$'' means $A\to B$, so $B$ is \emph{necessary} for $A$. ``$A$ if $B$'' means $B\to A$, so $B$ is \emph{sufficient} for $A$. ``If and only if'' asserts both directions. The consequence for optimisation problems is a discipline: a claim that the answer is $N$ decomposes into a \emph{sufficiency} half, an explicit procedure using $N$ steps, and a \emph{necessity} half, an argument that no procedure using $N-1$ steps can work. Neither half implies the other and an answer with only one half is incomplete.

\begin{worked}
A polynomial $p$ has unknown degree and unknown coefficients, all of them non-negative integers. You may ask for the value $p(k)$ at any integer $k$ of your choosing, and you learn the answer before choosing the next $k$. How many questions do you need?

\emph{Sufficiency: two suffice.} Ask $p(1)$. Since the coefficients are non-negative integers, $p(1)=S$ is their sum, so every coefficient lies in $\{0,1,\ldots,S\}$. Now ask $p(S+1)$. Because each coefficient is strictly less than the base $S+1$, the value $p(S+1)=\sum_i c_i(S+1)^i$ is precisely the base-$(S+1)$ representation of that integer, and the coefficients are its digits. Reading the digits recovers $p$ completely. Random testing of this procedure on four thousand polynomials of degree up to six with coefficients up to forty recovers every one of them exactly, verified in \texttt{python3}.

\emph{Necessity: one does not suffice.} A single question is a single $k$ fixed in advance. Whatever $k$ you choose, the two distinct polynomials $p(x)=x^2$ and $q(x)=kx$ both have non-negative integer coefficients and both return $k^2$, so no single value can separate them. Answer: exactly two, and note that the two halves used completely different ideas, the first a construction and the second a collision.
\end{worked}

\begin{trap}
Producing a clever procedure and calling it optimal. A construction proves an upper bound on the cost, never a lower one, and the sentence ``I cannot see how to do better'' is not an argument. The opposite failure is equally common in written solutions: a slick counting bound with no strategy exhibited, so that the claimed minimum may not be achievable at all. Write both halves down separately and label them, because a solution that silently merges them usually contains only one.
\end{trap}

\begin{seealsob}Contrapositive, converse, inverse; the extremal principle; the material conditional and vacuous truth.\end{seealsob}

\begin{keynote}[Translating the six phrases]
Six English constructions cover most of what puzzle statements do with conditionals, and it is worth having them memorised as a table rather than reasoning them out each time. ``$A$ implies $B$'', ``$A$ only if $B$'', ``$B$ whenever $A$'', ``$A$ is sufficient for $B$'' and ``$B$ is necessary for $A$'' all mean $A\to B$. ``$A$ if $B$'', ``$B$ implies $A$'' and ``$A$ is necessary for $B$'' all mean $B\to A$. The two lists are exchanged by swapping $A$ and $B$, and the single phrase that catches people out is ``only if'', which points the arrow the opposite way from the way it reads. When a problem statement mixes several of these in one paragraph, translate every sentence into arrow form before doing anything else; the mixture is usually the difficulty.
\end{keynote}

\section{Recognition matrix}
\begin{recog}{@{}p{0.30\textwidth}p{0.36\textwidth}p{0.28\textwidth}@{}}
\rowcolor{mist}\mh{If you see} & \mh{Reach for} & \mh{If that fails} \\
\midrule
\endhead
Three or four atoms, several constraints & Enumerate $2^n$ rows; keep the survivors & Case on the most constrained atom \\
``Does this follow?'' & Entailment: check \emph{every} survivor & One counterexample row refutes it \\
``Is this possible?'' & Satisfiability: one survivor suffices & Zero survivors means impossible \\
A unique survivor & Fully determined; report it & Several survivors: report only the shared part \\
Information given about $B$, rule is $A\to B$ & Contrapositive $\lnot B\to\lnot A$ & You do not own the converse \\
Asked which case to test & The unique false row: $A$ true, $B$ false & Ignore rows that can only confirm \\
``Not both'' & $\lnot A\vee\lnot B$: two branches & Not ``neither'', which is $\lnot A\wedge\lnot B$ \\
``None of them'' & $\lnot A_1\wedge\cdots\wedge\lnot A_n$ & A list of facts, not a case split \\
``At most one'' & $\bigwedge_{i<j}(\lnot A_i\vee\lnot A_j)$ & Add a disjunction for ``exactly one'' \\
A negated compound & Push $\lnot$ inward to negation normal form & Then read the top-level connective \\
Refute ``for all'' & Exhibit one witness with $\lnot P$ & Search a small finite case first \\
Refute ``there exists'' & Prove $\lnot P$ over the whole domain & An invariant or a parity argument \\
Nested $\forall\exists$ & Negate outside-in, keep the order & Ask whether $y$ may depend on $x$ \\
A conditional with a doubtful antecedent & $\lnot A\vee B$; check the vacuous branch & The rule is free when $A$ fails \\
``If I am a knight then \ldots'' & $k\leftrightarrow(k\to S)$ forces $k$ and $S$ & The knave version determines nothing \\
A rule over a possibly empty set & Vacuously true; attack the emptiness & Show the range is nonempty first \\
``Only if'' & $A\to B$, arrow away from the ``only'' & Reverses the naive reading \\
``If'' in a permission or a promise & Suspect a hidden biconditional & State which reading you adopted \\
``Minimum number of \ldots'' & Construction plus lower bound, separately & One half alone is not an answer \\
A clever strategy, no bound & You have sufficiency only & Find a collision of two inputs \\
\end{recog}
